\documentclass{beamer}
\usepackage{bussproofs}
\usepackage{graphicx}
\usepackage{multicol}

\title{Curry-Howard Correspondence}
\author{Tristan McDermott}
\institute{Advisor: Tibor Beke}
\date{December 12, 2025}

\begin{document}

\frame{\titlepage}

% give them a "soft start", almost like history: lambda calc was invented by church as a foundation of math, so was intuitionistic logic (they thought proofs by contradiction were stupid), then much later curry and howard noticed the correspondence
\begin{frame}{Introduction}
    \begin{itemize}
        \item \textbf{Curry-Howard correspondence:} establishes a connection between logic and type theory
        \item ``Proofs are programs'' (and vice versa)
        \item First noticed by Haskell Curry in 1958 (but between different systems we won't talk about)
        \item Added onto by William Howard in 1969
        \item He noticed a connection between \textbf{lambda calculus} and \textbf{natural deduction}
    \end{itemize}
\end{frame}

\begin{frame}{Lambda calculus}
    \begin{itemize}
        \item \textbf{Lambda calculus:} a formal language dealing with functions
        \item Basically, a very simple programming language
        \item Terms have 3 forms:
        \begin{itemize}
            \item Variable: $x$
            \item Abstraction: $\lambda x. M$
            \item Application: $M N$
        \end{itemize}
        \item $M$ and $N$ are themselves terms
        \item An abstraction creates a function, an application calls it
        % bound variables?
    \end{itemize}
\end{frame}

\begin{frame}{$\beta$-reduction}
    \begin{itemize}
        \item Lambda calculus terms can be simplified by \textbf{$\beta$-reduction}
        \item $(\lambda x. M) N \to_\beta M[N/x]$
        \item Substitution only replaces free instances of $x$
        \item \textbf{$\beta$-normal form:} the form a term has when it can't be $\beta$-reduced anymore
        \item \textbf{Church-Rosser property:} guarantees all terms which have a $\beta$-normal form have a unique $\beta$-normal form
        \item Are there terms without $\beta$-normal forms at all?
    \end{itemize}
\end{frame}

\begin{frame}{Infinite $\beta$-reduction}
    \begin{itemize}
        \item $(\lambda x. xx)(\lambda x. xx) \newline
        \to_\beta (\lambda x. xx)(\lambda x. xx) \newline
        \to_\beta (\lambda x. xx)(\lambda x. xx) \newline
        \to_\beta ...$
        \item If the lambda calculus is a programming language, this means some programs never terminate!
        \item Problem: any function can take any kind of input
        \item How can we prevent this?
    \end{itemize}
\end{frame}

\begin{frame}{Simply-typed lambda calculus}
    \begin{itemize}
        \item Add types to untyped lambda calculus
        \item \textbf{Type:} hard to define, but every term has exactly 1 type
        \begin{itemize}
            \item Basic types: $0$ (no terms) and $1$ (1 term: $*$)
            \item Function type: $A \to B$ where $A$ and $B$ are types
            \item $|A \to B| = |B|^{|A|}$
        \end{itemize}
        \item Untyped lambda calculus + the types shown above = \textbf{simply-typed lambda calculus}
    \end{itemize}
\end{frame}

\begin{frame}{Simply-typed lambda calculus}
    \begin{itemize}
        \item Abstractions look different: $\lambda x^A.M$ where $A$ is $x$'s type
        \item Application arguments much match the abstraction's parameter's type
        \item \textbf{Strong normalization:} every term now has a $\beta$-normal form
        \item $(\lambda x^{1 \to 1}. x*)*$ is invalid
        \item $(\lambda x^{1 \to 1}. x*)(\lambda y^1. y)$ is valid
        \item How can we determine what terms are well-typed?
    \end{itemize}
\end{frame}

\begin{frame}{Typing rules}
    \begin{prooftree}
        \AxiomC{}
        \RightLabel{$Ax$}
        \UnaryInfC{$\Gamma, x: A \vdash x: A$}
    \end{prooftree}
    
    \begin{prooftree}
        \AxiomC{$\Gamma, x: A \vdash M: B$}
        \RightLabel{$\to I$}
        \UnaryInfC{$\Gamma \vdash \lambda x^A. M: A \to B$}
    \end{prooftree}
    
    \begin{prooftree}
        \AxiomC{$\Gamma \vdash M: A \to B$}
        \AxiomC{$\Gamma \vdash N: A$}
        \RightLabel{$\to E$}
        \BinaryInfC{$\Gamma \vdash MN: B$}
    \end{prooftree}
    
    \begin{prooftree}
        \AxiomC{}
        \RightLabel{$1I$}
        \UnaryInfC{$\Gamma \vdash *: 1$}
    \end{prooftree}
    
    \begin{prooftree}
        \AxiomC{$\Gamma \vdash M: 0$}
        \RightLabel{$0E$}
        \UnaryInfC{$\Gamma \vdash \text{abort}_A M: A$}
    \end{prooftree}
\end{frame}

% maybe instead of "constructive" say "direct, not indirect reasoning"
\begin{frame}{Natural deduction}
    \begin{itemize}
        \item \textbf{Natural deduction:} a proof calculus based on inference rules rather than axioms
        \item \textbf{Proposition:} a statement with a truth value
        \begin{itemize}
            \item Basic proposition: $A$
            \item Implication: $A \implies B$ where $A$ and $B$ are propositions
        \end{itemize}
        \item We will use NJ, which is based on \textbf{intuitionistic logic}
        \item No law of excluded middle or double negation elimination
        \item Meant to encourage constructive proofs
    \end{itemize}
\end{frame}

\begin{frame}{Inference rules}
    \begin{prooftree}
        \AxiomC{}
        \RightLabel{$Ax$}
        \UnaryInfC{$\Gamma, x: A \vdash x: A$}
    \end{prooftree}
    
    \begin{prooftree}
        \AxiomC{$\Gamma \vdash B$}
        \RightLabel{$\implies I_v$}
        \UnaryInfC{$\Gamma \vdash A \implies B$}
    \end{prooftree}
    
    \begin{prooftree}
        \AxiomC{$\Gamma, x: A \vdash B$}
        \LeftLabel{$x$}
        \RightLabel{$\implies I_s$}
        \UnaryInfC{$\Gamma \vdash A \implies B$}
    \end{prooftree}
    
    \begin{prooftree}
        \AxiomC{$\Gamma \vdash A \implies B$}
        \AxiomC{$\Gamma \vdash A$}
        \RightLabel{$\implies E$}
        \BinaryInfC{$\Gamma \vdash B$}
    \end{prooftree}
    
    \begin{prooftree}
        \AxiomC{}
        \RightLabel{$\top I$}
        \UnaryInfC{$\Gamma \vdash \top$}
    \end{prooftree}
    
    \begin{prooftree}
        \AxiomC{$\Gamma \vdash \bot$}
        \RightLabel{$\bot E$}
        \UnaryInfC{$\Gamma \vdash A$}
    \end{prooftree}
\end{frame}

\begin{frame}{Detour elimination}
    \begin{itemize}
        \item \textbf{Detour:} an introduction rule followed immediately by an elimination rule
        \begin{prooftree}
            \AxiomC{$\Gamma \vdash B$}
            \RightLabel{$\implies I_v$}
            \UnaryInfC{$\Gamma \vdash \top \implies B$}
            \AxiomC{}
            \RightLabel{$\top I$}
            \UnaryInfC{$\Gamma \vdash \top$}
            \RightLabel{$\implies E$}
            \BinaryInfC{$\Gamma \vdash B$}
        \end{prooftree}
        \item All this proves is B, given B
        \item A proof with no detours is \textbf{normal}
        \item \textbf{Confluence:} every proof that has a normal form has a unique normal form
        \item \textbf{Strong normalization:} every proof has a normal form
    \end{itemize}
\end{frame}

\begin{frame}{The correspondence}
    \renewcommand{\arraystretch}{2}
    \resizebox{\linewidth}{!}{
        \begin{tabular}{c c}
            \begin{minipage}{0.55\textwidth}
                \begin{prooftree}
                    \AxiomC{}
                    \RightLabel{$Ax$}
                    \UnaryInfC{$\Gamma, x: A \vdash x: A$}
                \end{prooftree}
            \end{minipage} & \begin{minipage}{0.55\textwidth}
                \begin{prooftree}
                    \AxiomC{}
                    \RightLabel{$Ax$}
                    \UnaryInfC{$\Gamma, x: A \vdash x: A$}
                \end{prooftree}
            \end{minipage} \\
            What goes here? & \begin{minipage}{0.55\textwidth}
                \begin{prooftree}
                    \AxiomC{$\Gamma \vdash B$}
                    \RightLabel{$\implies I_v$}
                    \UnaryInfC{$\Gamma \vdash A \implies B$}
                \end{prooftree}
            \end{minipage} \\
            \begin{minipage}{0.55\textwidth}
                \begin{prooftree}
                    \AxiomC{$\Gamma, x: A \vdash M: B$}
                    \RightLabel{$\to I$}
                    \UnaryInfC{$\Gamma \vdash \lambda x^A. M: A \to B$}
                \end{prooftree}
            \end{minipage} & \begin{minipage}{0.55\textwidth}
                \begin{prooftree}
                    \AxiomC{$\Gamma, x: A \vdash B$}
                    \LeftLabel{$x$}
                    \RightLabel{$\implies I_s$}
                    \UnaryInfC{$\Gamma \vdash A \implies B$}
                \end{prooftree}
            \end{minipage} \\
            \begin{minipage}{0.55\textwidth}
                \begin{prooftree}
                    \AxiomC{$\Gamma \vdash M: A \to B$}
                    \AxiomC{$\Gamma \vdash N: A$}
                    \RightLabel{$\to E$}
                    \BinaryInfC{$\Gamma \vdash MN: B$}
                \end{prooftree}
            \end{minipage} & \begin{minipage}{0.55\textwidth}
                \begin{prooftree}
                    \AxiomC{$\Gamma \vdash A \implies B$}
                    \AxiomC{$\Gamma \vdash A$}
                    \RightLabel{$\implies E$}
                    \BinaryInfC{$\Gamma \vdash B$}
                \end{prooftree}
            \end{minipage} \\
            \begin{minipage}{0.55\textwidth}
                \begin{prooftree}
                    \AxiomC{}
                    \RightLabel{$1I$}
                    \UnaryInfC{$\Gamma \vdash *: 1$}
                \end{prooftree}
            \end{minipage} & \begin{minipage}{0.55\textwidth}
                \begin{prooftree}
                    \AxiomC{}
                    \RightLabel{$\top I$}
                    \UnaryInfC{$\Gamma \vdash \top$}
                \end{prooftree}
            \end{minipage} \\
            \begin{minipage}{0.55\textwidth}
                \begin{prooftree}
                    \AxiomC{$\Gamma \vdash M: 0$}
                    \RightLabel{$0E$}
                    \UnaryInfC{$\Gamma \vdash \text{abort}_A M: A$}
                \end{prooftree}
            \end{minipage} & \begin{minipage}{0.55\textwidth}
                \begin{prooftree}
                    \AxiomC{$\Gamma \vdash \bot$}
                    \RightLabel{$\bot E$}
                    \UnaryInfC{$\Gamma \vdash A$}
                \end{prooftree}
            \end{minipage}
        \end{tabular}
    }
\end{frame}

\begin{frame}{The correspondence}
    \begin{prooftree}
        \AxiomC{}
        \RightLabel{$Ax$}
        \UnaryInfC{$\Gamma, x: \top \implies A \vdash \top \implies A$}
        \AxiomC{}
        \RightLabel{$\top I$}
        \UnaryInfC{$\Gamma \vdash \top$}
        \RightLabel{$\implies E$}
        \BinaryInfC{$\Gamma, x: \top \implies A \vdash A$}
        \LeftLabel{$x$}
        \RightLabel{$\implies I_s$}
        \UnaryInfC{$\Gamma \vdash (\top \implies A) \implies A$}
    \end{prooftree}
    
    \begin{center}
        \Huge{$\Updownarrow$}
    \end{center}
    
    \begin{prooftree}
        \AxiomC{}
        \RightLabel{$Ax$}
        \UnaryInfC{$\Gamma, x: 1 \to A \vdash x: 1 \to A$}
        \AxiomC{}
        \RightLabel{$1 I$}
        \UnaryInfC{$\Gamma \vdash *: 1$}
        \RightLabel{$\to E$}
        \BinaryInfC{$\Gamma, x: 1 \to A \vdash x *: A$}
        \RightLabel{$\to I$}
        \UnaryInfC{$\Gamma \vdash \lambda x^{1 \to A}.x *: (1 \to A) \to A$}
    \end{prooftree}
\end{frame}

\begin{frame}{Product types}
    \begin{itemize}
        \item \textbf{Product type:} $A \times B$ where $A$ and $B$ are types
        \item A pair of 2 terms, with the ability to extract each part from the whole
        \item $|A \times B| = |A| \times |B|$
    \end{itemize}
    
    \begin{prooftree}
        \AxiomC{$\Gamma \vdash M: A$}
        \AxiomC{$\Gamma \vdash N: B$}
        \RightLabel{$\times I$}
        \BinaryInfC{$\Gamma \vdash \langle M, N \rangle: A \times B$}
    \end{prooftree}
    
    \begin{prooftree}
        \AxiomC{$\Gamma \vdash M: A \times B$}
        \RightLabel{$\times E_1$}
        \UnaryInfC{$\Gamma \vdash \pi_1 M: A$}
    \end{prooftree}
    
    \begin{prooftree}
        \AxiomC{$\Gamma \vdash M: A \times B$}
        \RightLabel{$\times E_2$}
        \UnaryInfC{$\Gamma \vdash \pi_2 M: B$}
    \end{prooftree}
\end{frame}

\begin{frame}{Conjunctions}
    \begin{itemize}
        \item \textbf{Conjunction:} $A \land B$ where $A$ and $B$ are propositions
        \item ``Logical AND''
        \item Both propositions are true, and each one is true separately
    \end{itemize}

    \begin{prooftree}
        \AxiomC{$\Gamma \vdash A$}
        \AxiomC{$\Gamma \vdash B$}
        \RightLabel{$\land I$}
        \BinaryInfC{$\Gamma \vdash A \land B$}
    \end{prooftree}

    \begin{prooftree}
        \AxiomC{$\Gamma \vdash A \land B$}
        \RightLabel{$\land E_1$}
        \UnaryInfC{$\Gamma \vdash A$}
    \end{prooftree}
    
    \begin{prooftree}
        \AxiomC{$\Gamma \vdash A \land B$}
        \RightLabel{$\land E_2$}
        \UnaryInfC{$\Gamma \vdash B$}
    \end{prooftree}
\end{frame}

\begin{frame}{Product types and conjunction}
    \begin{itemize}
        \item Both are combinations of 2 things, and you can ``extract'' each part
        \item Conjunctions with false are all false: $A \land \bot \iff \bot$...
        \item and products with 0 don't exist: $A \times 0 = 0$
        \item Conjunctions with true do nothing: $A \land \top \iff A$...
        \item and products with 1 are pointless: $A \times 1 = A$
    \end{itemize}
\end{frame}

% mention its a disjoint sum
\begin{frame}{Sum types}
    \begin{itemize}
        \item \textbf{Sum type:} $A + B$ where $A$ and $B$ are types
        \item Either an $A$ or a $B$, with a ``tag'' that tells you which one it is
        \item $|A + B| = |A| + |B|$
    \end{itemize}

    \begin{prooftree}
        \AxiomC{$\Gamma \vdash M: A$}
        \RightLabel{$+ I_1$}
        \UnaryInfC{$\Gamma \vdash \text{inj}_1 M: A + B$}
    \end{prooftree}
    
    \begin{prooftree}
        \AxiomC{$\Gamma \vdash M: B$}
        \RightLabel{$+ I_2$}
        \UnaryInfC{$\Gamma \vdash \text{inj}_2 M: A + B$}
    \end{prooftree}
    
    \begin{prooftree}
        \AxiomC{$\Gamma \vdash M: A + B$}
        \AxiomC{$\Gamma, x: A \vdash N: C$}
        \AxiomC{$\Gamma, y: B \vdash O: C$}
        \RightLabel{$+ E$}
        \TrinaryInfC{$\Gamma \vdash \text{case } M \text{ of } (x)N; (y)O: C$}
    \end{prooftree}
\end{frame}

\begin{frame}{Disjunction}
    \begin{itemize}
        \item \textbf{Disjunction:} $A \lor B$ where $A$ and $B$ are propositions
        \item ``Logical OR''
        \item At least 1 proposition is true, but we can't always tell which one
    \end{itemize}

    \begin{prooftree}
        \AxiomC{$\Gamma \vdash A$}
        \RightLabel{$\lor I_1$}
        \UnaryInfC{$\Gamma \vdash A \lor B$}
    \end{prooftree}
    
    \begin{prooftree}
        \AxiomC{$\Gamma \vdash B$}
        \RightLabel{$\lor I_2$}
        \UnaryInfC{$\Gamma \vdash A \lor B$}
    \end{prooftree}
    
    \begin{prooftree}
        \AxiomC{$\Gamma \vdash A \lor B$}
        \AxiomC{$\Gamma, x: A \vdash C$}
        \AxiomC{$\Gamma, y: B \vdash C$}
        \LeftLabel{$x, y$}
        \RightLabel{$\lor E$}
        \TrinaryInfC{$\Gamma \vdash C$}
    \end{prooftree}
\end{frame}

% add parentheses in distributivty
\begin{frame}{Sum types and disjunction}
    \begin{itemize}
        \item Commutative: $A + B = B + A$ and $A \lor B \iff B \lor A$
        \item Associative: $(A + B) + C = A + (B + C)$ and $(A \lor B) \lor C \iff A \lor (B \lor C)$
        \item Identity element: $A + 0 = A$ and $A \lor \bot \iff A$
        \item Distributive: $A \times (B + C) = A \times B + A \times C$ and $A \land (B \lor C) \iff A \land B \lor A \land C$
        \item Dual to products and conjunctions
    \end{itemize}
\end{frame}

% is it really "2nd order NJ?"
\begin{frame}{Limitations of STLC and NJ}
    \begin{itemize}
        \item Can't express universals
        \item No single function $\lambda x^A.x$ exists for all types $A$
        \item $A \implies A$, $B \implies B$, etc. all have to be proven separately
        \item STLC + polymorphism = \textbf{System F}
        \item NJ + universal quantification = \textbf{2nd-order NJ}
    \end{itemize}
\end{frame}

\begin{frame}{System F}
    \begin{itemize}
        \item \textbf{Universal type:} $\forall \alpha. A$ where $A$ is a type
        \item Functions that take a type and return a term: $\Lambda \alpha. \lambda x^\alpha. x: \forall \alpha. \alpha \to \alpha$
        \item New $\beta$-reduction rule: $(\Lambda \alpha. A) B \to_\beta A[B/\alpha]$
        \item Still Church-Rosser and strongly normalizing
    \end{itemize}
    
    \begin{prooftree}
        \AxiomC{$\Gamma \vdash M: A$}
        \AxiomC{$\alpha \not\in FTV(\Gamma)$}
        \RightLabel{$\forall I$}
        \BinaryInfC{$\Gamma \vdash \Lambda \alpha. M: \forall \alpha. A$}
    \end{prooftree}
    
    \begin{prooftree}
        \AxiomC{$\Gamma \vdash \Lambda \alpha. M: \forall \alpha. A$}
        \RightLabel{$\forall E$}
        \UnaryInfC{$\Gamma \vdash M: A[B/\alpha]$}
    \end{prooftree}
\end{frame}

\begin{frame}{2nd-order NJ}
    \begin{itemize}
        \item \textbf{Universal quantifier:} $\forall \alpha. A$ where $A$ is a proposition
        \item ``For all $\alpha$, $A$''
    \end{itemize}

    \begin{prooftree}
        \AxiomC{$\Gamma \vdash A$}
        \AxiomC{$\alpha \not\in FV(\Gamma)$}
        \RightLabel{$\forall I$}
        \BinaryInfC{$\Gamma \vdash \forall \alpha. A$}
    \end{prooftree}
    
    \begin{prooftree}
        \AxiomC{$\Gamma \vdash \forall \alpha.A$}
        \RightLabel{$\forall E$}
        \UnaryInfC{$\Gamma \vdash A[B/\alpha]$}
    \end{prooftree}
    
    \begin{itemize}
        \item New introduction and elimination rule means new type of detour:
    \end{itemize}

    \begin{prooftree}
        \AxiomC{$\Gamma \vdash A$}
        \AxiomC{$\alpha \not\in FV(\Gamma)$}
        \RightLabel{$\forall I$}
        \BinaryInfC{$\Gamma \vdash \forall \alpha. A$}
        \RightLabel{$\forall E$}
        \UnaryInfC{$\Gamma \vdash A[B/\alpha]$}
    \end{prooftree}
\end{frame}

\begin{frame}{The identities}
    \begin{itemize}
        \item The connectives follow 14 identities:
    \end{itemize}

    \begin{multicols}{2}
        \begin{enumerate}
            \item $0 + x = x$
            \item $x + y = y + x$
            \item $(x + y) + z = x + (y + z)$
            \item $x \times 0 = 0$
            \item $x \times 1 = x$
            \item $x \times y = y \times x$
            \item $(x \times y) \times z = x \times (y \times z)$
            \item $x \times (y + z) = x \times y + x \times z$
            \item $0 \to x = 1$
            \item $x \to 1 = 1$
            \item $1 \to x = x$
            \item $x + y \to z = (x \to z) \times (y \to z)$
            \item $x \to y \times z = (x \to y) \times (x \to z)$
            \item $x \to y \to z = x \times y \to z$
        \end{enumerate}
    \end{multicols}
    
    \begin{itemize}
        \item (Keep in mind $|A \to B| = |B|^{|A|}$)
        \item What algebraic structure do these operations and axioms create?
    \end{itemize}
\end{frame}

% dont worry too much about exponential/product/coproduct, just say they behave like functions, cartesian products, and disjoint union
\begin{frame}{Categories}
    \begin{itemize}
        \item A Bicartesian closed category... what is that?
        \item \textbf{Category:} a collection of objects and morphisms
        \item They can be many things:
        \begin{itemize}
            \item Sets and functions
            \item Vector spaces and linear transformations
            \item Groups and group homomorphisms
        \end{itemize}
        \item Morphisms must be composable: $f: A \to B$ and $g: B \to C$ means there must be $g \circ f: A \to C$
        \item Every object needs an identity morphism $id_A: A \to A$ which acts as an identity for composition
    \end{itemize}
\end{frame}

% explain what a BCC is here: it has exponential, product, and sum objects for all pairs of objects, and an initial and terminal object
% include NOTHING from theorem proving section, maybe mention system f in the context of software engineering and how lambda calc based type systems make functional languages special
\begin{frame}{Lambek}
    \begin{itemize}
        \item Category theoretic contributions made by Joachim Lambek
        \item True name of the correspondence: \textbf{Curry-Howard-Lambek}
    \end{itemize}
    
    \renewcommand{\arraystretch}{1.5}
    \resizebox{\linewidth}{!}{
        \begin{tabular}{c c c c c}
            STLC & $\iff$ & NJ & $\iff$ & Bicartesian closed category \\
            Term & $\iff$ & Proof & $\iff$ & Morphism \\
            Type & $\iff$ & Proposition & $\iff$ & Object \\
            Typing rule & $\iff$ & Inference rule & $\iff$ & Categorical axiom \\
            $\beta$-reduction & $\iff$ & Detour elimination & $\iff$ & Morphism composition \\
            1 & $\iff$ & $\top$ & $\iff$ & Terminal object \\
            0 & $\iff$ & $\bot$ & $\iff$ & Initial object \\
            Function type & $\iff$ & Implication & $\iff$ & Exponential object \\
            Product type & $\iff$ & Conjunction & $\iff$ & Product \\
            Sum type & $\iff$ & Disjunction & $\iff$ & Coproduct
        \end{tabular}
    }
\end{frame}

\begin{frame}{Lambek}
    \begin{itemize}
        \item Why is Lambek left out of most discussions of the correspondence?
        \begin{itemize}
            \item Logic lets you prove theorems
            \item Type systems let you write safe programs
            \item What can categories do for the average mathematician?
        \end{itemize}
        \item Why is it useful despite that?
        \begin{itemize}
            \item STLC and NJ are pure syntax
            \item BCCs provide semantics as mathematical objects that model them
            \item Category theory lets us reason about these systems more abstractly and generally
        \end{itemize}
    \end{itemize}
\end{frame}

\begin{frame}{}
    \begin{center}
        \Huge{Questions?}
    \end{center}
\end{frame}

\begin{frame}{}
    \begin{center}
        \Huge{Thank you}
    \end{center}
\end{frame}

\end{document}
